\section{A polynomial improvement by scheduling adjacent blocks}
\label{sec:polynomial-blocks}

The individual-vertex schedule of Section~\ref{sec:mass-improvement}
admits a stronger polynomial-time version. We apply the same exact-
coverage and geometric-batch algorithm to adjacent hull edges, and
account for discovery and endpoint decoding in one cost bound.
This first gives a factor below $5.312551$. Retaining orientation
credits from the first two batches further lowers the factor to
$5.283310$ for arbitrary-prior expected depth in the general
supplied-line model. Later refinements change the rational schedule
score and its certificate, while reusing the stage-to-tree conversion.
The proof has one ledger: discovery gives an optimal-path restriction
credit, endpoint recovery leaves a conditional-entropy cost, and
additional block savings reduce the remaining coefficient of positive
mass.

\subsection{The polynomial block algorithm}

Use the mass-balanced and mass-flat recursion already described, with
the new rational threshold
\[
 \delta=\frac{138821934}{10^9},\qquad
 \beta=\frac{3512863}{10^6}.
\]
In a flat state, keep the full normalized hull, including the origin
if it is a hull vertex. Its edges are candidate blocks. A block queries
its two real endpoints in a fixed order, omitting cached answers; an
origin endpoint is a free constant-negative marker. Both endpoint tests
are supplied lines. Retain the origin in the cyclic hull order so every
candidate block remains an edge of the full hull.

Every nonzero positive-vertex interval meets a circular interval of
edges. These are the block-hit demands. Apply
Lemma~\ref{lem:mssc} to this incidence system, choosing the shift by
the rational objective
\[
 \MSSC_{\mathrm{blk}}(G)+p_0|G|,
 \qquad p_0=w(A_0),\quad b=w(S)-p_0.
\]
Here $p_0=0$ when the zero pattern is absent. Query blocks until a
positive endpoint answer or an earlier singleton inference; then learn
the full hull pattern with the existing alphabetic endpoint codes and
recurse on its class. A completed negative cover identifies the zero
pattern, when present.

Let $k_A$ be the first block hitting a nonzero class $A$, and let
$K=|G|$. If the first queried endpoint is positive, discovery costs
at most $2k_A-1$ and the endpoint codes have overhead at most four.
If the second endpoint is the first positive answer, discovery costs
at most $2k_A$ and its adjacent negative endpoint pins one boundary,
leaving overhead at most two. An origin block has only one real query.
Thus every positive route has nonentropy allowance at most
\[
 2(k_A-1)+5=2k_A+3.
\]
For the actual discovery partition $\mathcal R$, write
$J=I(S,\mathcal R)$ and $I=I(S,\Pi)$, where $\Pi$ is the full-pattern
partition. The chain-rule proof of Lemma~\ref{lem:pruning-macro} gives
\begin{equation}\label{eq:poly-block-stage}
 W_{\mathrm{pat}}(S)\le
 2\MSSC_{\mathrm{blk}}(G)+3b+2p_0K+\frac{I-J}{\ln2}.
\end{equation}
The coefficient is $3b$ for ordinary first-hit ranks, or $5b$ for ranks
shifted by one. Using $5b$ with ordinary ranks would lose the saving.

Each nonzero discovery exit is contained in the cap of its actual
first positive endpoint, whose mass is less than $\delta w(S)$.
Although a block's union can have mass up to $2\delta w(S)$,
Lemma~\ref{lem:capped-discovery-entropy} still applies at the smaller
individual cap bound:
\begin{equation}\label{eq:poly-block-entropy}
 J\ge b\sigma(\delta),\qquad
 \sigma(\delta)=-k\delta\ln\delta-r\ln r,
 \quad k=\lfloor1/\delta\rfloor,\ r=1-k\delta.
\end{equation}
Early singleton exits only improve this bound and the query cost.

\subsection{Restriction credit and discrete block savings}

Keep the functions $\Phi_\beta$, $Z_\beta$, $\lambda_\beta$,
$\nu_\beta$, and $\eta_\beta$ from Lemma~\ref{lem:mssc}. Put
\[
 e_\beta(t)=2\bigl[\lambda_\beta t+\eta_\beta-\Phi_\beta(t)\bigr].
\]
Then $e_\beta(1)=0$, and
\[
 e_\beta(t)=\frac1{\ln\beta}\int_1^t
 \left[\beta+1-\frac{\lfloor x\rfloor+\lfloor\beta x\rfloor}{x}\right]dx.
\]
The integrand is nonnegative, so these discounts are nondecreasing.
In particular,
\begin{equation}\label{eq:poly-block-phi}
 2\Phi_\beta(t)=2\lambda_\beta(t-1)+2\Phi_\beta(1)-e_\beta(t).
\end{equation}

Follow a common witness through an expected-optimal tree, using the
$q$ queries and supporting-hit indices $t_A$ of
Lemma~\ref{lem:support-prefix}. Replace each witness query by a full-hull
edge containing its one or two supports; a singleton support may use
either incident edge. Enlarging a support set can only advance a hit.
Delete repeated and unproductive comparison blocks. If the zero
pattern is absent and the terminal singleton remains unhit, append an
edge incident to one of its positive vertices. Its rank is at most
$q+1=t_A$. All other classes are hit by rank at most $t_A$ too.
If $A_0$ is present, no append is needed and the minimum block cover
size $k_0$ is at most $q$.

Write
\[
 D_+=\sum_{A\ne A_0}w(A)(t_A-1),\qquad
 \Delta=\Delta_{\mathrm E}(S,\Pi).
\]
Supporting-prefix credit gives $D_++p_0q\le\Delta$.
Monotonicity of $\Phi_\beta$ and Lemma~\ref{lem:mssc} give
\[
 \MSSC_{\mathrm{blk}}(G)+p_0K
 \le\sum_{A\ne A_0}w(A)\Phi_\beta(t_A)+p_0 Z_\beta(k_0).
\]
The zero term vanishes when $p_0=0$, so no cover-size comparison is
needed in that case. For the chosen $\beta$, $\lambda_\beta>\nu_\beta$.
Combining \eqref{eq:poly-block-stage}, \eqref{eq:poly-block-phi}, and
$Z_\beta(k_0)\le\nu_\beta k_0-1/\ln\beta$ yields
\begin{equation}\label{eq:poly-block-discrete}
 \begin{aligned}
 W_{\mathrm{pat}}(S)\le{}&
 A\Delta+[3+2\Phi_\beta(1)]b
 -\sum_t m_t e_\beta(t)-\frac{2p_0}{\ln\beta}
 +\frac{I-J}{\ln2},\\
 A={}&2\lambda_\beta=\frac{\beta+1}{\ln\beta}.
 \end{aligned}
\end{equation}
Here $m_t$ is the mass first support-hit at witness index $t$, including
an appended terminal singleton. Each such analytical support block is
contained in at most two hull caps, hence $m_t<2\delta w(S)$.
These witness masses are separate from the actual discovery exits in
\eqref{eq:poly-block-entropy}; no equality between their ranks or
partitions is assumed.

\begin{center}
\small
\begin{tabular}{@{}p{.25\linewidth}p{.34\linewidth}p{.34\linewidth}@{}}
\hline
Object & Computed by the algorithm & Used for comparison only\\
\hline
First-hit rank & Block rank $k_A$ in the chosen schedule $G$
 & Supporting-prefix rank $t_A$ on an optimal witness path\\
Mass partition & Actual discovery exits, each inside one cap
 & Witness masses $m_t$, each inside at most two caps\\
Accounting role & Cost $\MSSC_{\rm blk}(G)$ and information $J$
 & Restriction credit $\Delta$ and discounts $e_\beta(t)$\\
\hline
\end{tabular}
\end{center}
Exact coverage connects these columns: a budget-$k$ batch covers at
least the witness prefix mass $M_k=\sum_{t\le k}m_t$. It does not
identify the algorithm's ranks with the witness ranks. This same
connection will pay for all later block-mode and initial-batch savings.

Define
\[
 N=\lfloor1/(2\delta)\rfloor,\quad s=1-2\delta N,\qquad
 \chi=2\delta\sum_{t=1}^{N}e_\beta(t)+s e_\beta(N+1).
\]
Apply Lemma~\ref{lem:capped-rank-credit} with
$a_t=e_\beta(t)$, $\tau=e_\beta(N+1)$,
$E=\{1,\ldots,N\}$, $\mu=0$, and $\zeta=2/\ln\beta$.
Monotonicity of $e_\beta$ makes these the only possible coefficients
below $\tau$; the lemma's $\chi$ equals the expression above.
If
\begin{equation}\label{eq:poly-block-credit-condition}
 e_\beta(N+1)\le\chi+\frac2{\ln\beta},
\end{equation}
the positive-rank discounts and the zero-route discount together save
at least $\chi b$. This is the entire mass-cap step: no distribution
of witness ranks is used beyond $m_t\le2\delta w(S)$.
Put
\begin{equation}\label{eq:poly-block-B}
 B=3+2\Phi_\beta(1)-\chi.
\end{equation}
Equation~\eqref{eq:poly-block-discrete} becomes the stronger local bound
\begin{equation}\label{eq:poly-block-local}
 W_{\mathrm{pat}}(S)\le A\Delta+B b+\frac{I-J}{\ln2}.
\end{equation}

If the normalized hull has dimension at most one, query its at most
two real extreme tests to learn the full pattern directly. Its cost
is at most $2w(S)$. Since $A>2$ and $B>2$, this is at most
$A\Delta+B b$: when the zero class exists, the witness path has
$q\ge1$ at a nonsingleton state and gives $\Delta\ge p_0$;
when it is absent, $b=w(S)$. Thus this case obeys the same local bound,
with $J=I$.

\subsection{The general polynomial guarantee}

\begin{theorem}[Improved polynomial expected-depth approximation]
\label{thm:poly-block-expected}
For a separating supplied family of binary unit-cost planar line queries
and arbitrary strictly positive rational priors, the preceding
deterministic polynomial-time algorithm returns a tree $T_{\mathrm{blk}}$
such that
\[
 c_{\mathrm E}(T_{\mathrm{blk}}(S))
 \le A\OPT_{\mathrm E}(S)+c\mathcal E(S),\qquad
 c=\max\{1/h(\delta),\ B/\sigma(\delta)\}.
\]
For the rational parameters above,
\begin{equation}\label{eq:poly-block-factor}
 \boxed{c_{\mathrm E}(T_{\mathrm{blk}}(\HH))
 \le\rho_{\mathrm{blk}}\OPT_{\mathrm E}(\HH),\qquad
 \rho_{\mathrm{blk}}=A+c\ln2<5.312551.}
\end{equation}
The running time is polynomial in the geometric input and rational bit
length, with no dependence of its exponent on the smallest prior.
\end{theorem}
\begin{proof}
Apply Lemma~\ref{lem:stage-conversion} to
\eqref{eq:poly-block-local} and \eqref{eq:poly-block-entropy}.
This proves the additive guarantee and its stated ratio.

For completeness, the numerical certificate uses exact rational
logarithm enclosures obtained from the atanh series
\eqref{eq:rational-log-bounds}. Here $N=3$, $\lfloor1/\delta\rfloor=7$,
and
\[
 B=7-\frac{\ln6}{\ln\beta}
       -2\delta[e_\beta(2)+e_\beta(3)]
       -(1-6\delta)e_\beta(4).
\]
Those enclosures certify $\lambda_\beta>\nu_\beta$ and
\[
 \begin{aligned}
 3.591810178840038&<A<3.591810178840040,\\
 5.013509704362398&<B<5.013509704362401,\\
 e_\beta(4)&<\chi+2/\ln\beta-1.08407,\\
 \max\{\ln2/h(\delta),\ B\ln2/\sigma(\delta)\}
 &<1.720740274489660.
 \end{aligned}
\]
In particular \eqref{eq:poly-block-credit-condition} holds and
$A+c\ln2<5.312550453329700<5.312551$.
The standalone certificate \texttt{poly-expected-check.py} verifies
the displayed inequalities with exact rational interval bounds.

The block demands are circular intervals, so all exact budgeted
coverage and minimum-cover calls are polynomial. There are at most
as many demanded patterns as hypotheses; the minimum cover size is
no larger. For fixed $\beta$, the shift enumeration has polynomially
many breakpoints; rational input weights give rational rank-weight sums
for choosing its schedule.
Logarithms appear only in the analysis. The existing endpoint codes
are polynomial to construct. Every nonsingleton output class is a
proper subset; hence the recursive partition has $O(n)$ states and
the complete construction is polynomial.
\end{proof}

\subsection{Retaining the first-batch orientation credit}

In an actual scheduled block, let $x,y$ be the newly hit masses
exclusive to its endpoints and $z$ their shared mass. Querying the
smaller exclusive mass first makes the larger exclusive mass exit
second, with a pinned boundary. The nonentropy allowance improves
from $5(x+y+z)$ to $5(x+y+z)-\max\{x,y\}$, in addition to shifted
discovery cost. Each cap has mass below $\delta w(S)$, so
\[
 \max\{x,y\}\ge\max\{0,x+y+z-\delta w(S)\}.
\]
To turn this actual-schedule saving into a general guarantee, couple it
to the witness through the exact first-batch coverage oracle. Use
\[
 \delta_* =\frac{142275675}{10^9}
\]
as the balance threshold, keeping the same $\beta$. For each candidate
shifted schedule $G$, use the orientations just specified and compute
\[
 \operatorname{Pin}(G)=\sum_{\text{blocks of }G}\max\{x,y\}.
\]
Choose the shift minimizing the rational surrogate
\begin{equation}\label{eq:poly-oriented-surrogate}
 F(G)=2\MSSC_{\mathrm{blk}}(G)+2p_0|G|-\operatorname{Pin}(G).
\end{equation}
The query-cost bound is now $F(G)+3b+(I-J)/\ln2$.
Both orientations have the same block union and first-hit rank, so
all coverage sets and greedy block orders remain those of the original
polynomial batch construction. The actual transcript and its entropy
may change, but $J\ge b\sigma(\delta_*)$ holds for every shift and
orientation.

For a log-uniform shift $u\in[1,\beta)$, the first budget is
$k=\lfloor u\rfloor\in\{1,2,3\}$. Let $H$ be the positive mass covered
by that batch. It has at most $k$ productive blocks, with newly hit
masses $g_j$. Consequently
\[
 \operatorname{Pin}(G_u)\ge
 \sum_{j\text{ in first batch}}(g_j-\delta_* w(S))_+
 \ge (H-k\delta_*w(S))_+.
\]
Write $M_k=\sum_{t\le k}m_t$ for the analytical witness mass used
in \eqref{eq:poly-block-discrete}. A comparison edge cover hits this
mass within its first $k$ blocks. Exact at-most-$k$ coverage gives
$H\ge M_k$. If the first budget already reaches the minimum cover size,
the batch covers all positive mass and the same statement holds.
Thus, with
\[
 P_k=\frac{\ln(\min\{k+1,\beta\}/k)}{\ln\beta},\qquad
 \mu=\sum_{k=1}^{3}kP_k=3-\frac{\ln6}{\ln\beta},
\]
the same shift average used in Lemma~\ref{lem:mssc} satisfies
\begin{equation}\label{eq:poly-oriented-transfer}
 \mathbb E[\operatorname{Pin}(G_u)]
 \ge\sum_{k=1}^{3}P_k(M_k-k\delta_*w(S))
 =\sum_{t=1}^{3}m_t\sum_{k=t}^{3}P_k-\delta_*w(S)\mu.
\end{equation}
This is a transfer between actual coverage and witness mass; it counts
the orientation credit once. Minimizing
\eqref{eq:poly-oriented-surrogate} permits this negative term to be
retained under the same shift average as the first-cover and length
terms. Minimizing only $\MSSC+p_0|G|$ would not suffice.

Define the effective coefficients
\[
 a_t=e_\beta(t)+\sum_{k=t}^{3}P_k\quad(t\le3),\qquad
 a_t=e_\beta(t)\quad(t\ge4).
\]
Explicitly $a_1=1$, $a_2=e_\beta(2)+1-P_1$, and
$a_3=e_\beta(3)+P_3$. Exact logarithm bounds certify
$a_1,a_2,a_3<e_\beta(4)$. Monotonicity of $e_\beta$ gives
$a_t\ge e_\beta(4)$ for $t\ge4$. We impose no ordering on the first
three coefficients. Define
\[
 \begin{aligned}
 \chi_{\mathrm{or}}
 &=2\delta_*[e_\beta(2)+e_\beta(3)]
       +(1-6\delta_*)e_\beta(4)+\delta_*\mu,\\
 B_{\mathrm{or}}&=3+2\Phi_\beta(1)-\chi_{\mathrm{or}}.
 \end{aligned}
\]
These are the coefficients from Lemma~\ref{lem:capped-rank-credit}
with $\tau=e_\beta(4)$, $E=\{1,2,3\}$, the intercept $\mu$ in
\eqref{eq:poly-oriented-transfer}, and $\zeta=2/\ln\beta$.
Indeed $\sum_{t=1}^3\sum_{k=t}^3P_k=\sum_{k=1}^3kP_k=\mu$,
which explains the positive term $\delta_*\mu$ in $\chi_{\rm or}$.
The certified condition
$\chi_{\mathrm{or}}+2/\ln\beta\ge e_\beta(4)$ lets the lemma turn
all discounts into $\chi_{\rm or}b$. Minimizing the rational
surrogate retains them under the same shift average. Substitution in
\eqref{eq:poly-block-discrete} gives
\[
 W_{\mathrm{pat}}(S)\le A\Delta+B_{\mathrm{or}}b+(I-J)/\ln2.
\]
The degenerate-hull argument also applies because $B_{\mathrm{or}}>2$.

\begin{theorem}[Orientation-aware polynomial approximation]
\label{thm:poly-oriented-expected}
In the general binary supplied-line model with arbitrary positive
rational priors, the block algorithm with threshold $\delta_*$ and
shift objective \eqref{eq:poly-oriented-surrogate} has approximation
factor
\[
 \boxed{\rho_{\mathrm{or}}=
 A+\max\{\ln2/h(\delta_*),\ B_{\mathrm{or}}\ln2/\sigma(\delta_*)\}
 <5.286244.}
\]
Its running time is polynomial in the input and rational bit length.
\end{theorem}
\begin{proof}
Apply Lemma~\ref{lem:stage-conversion} to this local ledger and
\eqref{eq:poly-block-entropy}, with
\[
 c_{\mathrm{or}}=\max\{1/h(\delta_*),B_{\mathrm{or}}/\sigma(\delta_*)\}.
\]
The exact certificate \texttt{poly-expected-next-check.py} gives
\[
 \begin{aligned}
 0.771720051343&<\chi_{\mathrm{or}}<0.771720051345,\\
 4.802209647029&<B_{\mathrm{or}}<4.802209647032,\\
 \max\{\ln2/h(\delta_*),B_{\mathrm{or}}\ln2/\sigma(\delta_*)\}
 &<1.694433600947,\\
 e_\beta(4)-\max\{a_1,a_2,a_3\}&>0.06815.
 \end{aligned}
\]
It also certifies the zero-route condition. Hence
$\rho_{\mathrm{or}}<5.286243779788<5.286244$.
All values in \eqref{eq:poly-oriented-surrogate} are rational
rank-weight and exclusive-mass sums. For each enumerated shift the
surrogate and the best block orientations are computable by a
polynomial incidence scan; schedules and orientations are constant
within the same shift cells. The change preserves the
polynomial construction and has no minimum-prior dependence.
\end{proof}

\subsection{A mixed-prefix certificate below 5.283310}

The algorithm can be sharpened slightly further by changing its balance
threshold to
\[
 \delta_\dagger=\frac{142671046}{10^9}.
\]
Keep the same block orientations and shift objective
\eqref{eq:poly-oriented-surrogate}. Only the analysis mixes two lower
bounds on the same actual orientation credit; it does not add those
credits or change the polynomial scheduler.

For a shift $u$, put $k_0(u)=\lfloor u\rfloor$ and
$k_1(u)=\lfloor\beta u\rfloor$. The first two batches use at most
$k_0(u)+k_1(u)$ blocks and cover mass at least $M_{k_1(u)}$.
If either budget reaches the minimum cover size, the schedule finishes
and covers all positive mass, which only improves these bounds.
Thus $\operatorname{Pin}(G_u)$ is at least each of
\[
 M_{k_0(u)}-\delta_\dagger w(S)k_0(u),\qquad
 M_{k_1(u)}-\delta_\dagger w(S)[k_0(u)+k_1(u)].
\]
It is therefore at least their convex combination, with second-prefix
coefficient
\[
 \alpha(u)=
 \begin{cases}
 7963376/10^7,&1\le u<4/\beta,\\
 952407/10^6,&4/\beta\le u<5/\beta,\\
 1,&5/\beta\le u<6/\beta,\\
 0,&6/\beta\le u<\beta.
 \end{cases}
\]
On the three nonzero-coefficient ranges, $k_0(u)=1$ and $k_1(u)$ is
respectively 3, 4, or 5. Define their averaged second-prefix weights
\[
 q_0=\frac{7963376}{10^7}\frac{\ln(4/\beta)}{\ln\beta},\quad
 q_1=\frac{952407}{10^6}\frac{\ln(5/4)}{\ln\beta},\quad
 q_2=\frac{\ln(6/5)}{\ln\beta}.
\]
Relative to the first-prefix bound, each range adds $M_{k_1}-M_1$
and incurs an additional block-budget charge $k_1\delta_\dagger w(S)$.
Hence the combined witness discount and averaged pin credit obey
\[
 \sum_t m_t e_\beta(t)+\mathbb E[\operatorname{Pin}(G_u)]
 \ge\sum_t m_t\widetilde a_t-\delta_\dagger w(S)\widetilde\mu,
\]
where
\[
 \begin{aligned}
 \widetilde a_1&=1,\\
 \widetilde a_2&=e_\beta(2)+1-P_1+q_0+q_1+q_2,\\
 \widetilde a_3&=e_\beta(3)+P_3+q_0+q_1+q_2,\\
 \widetilde a_4&=e_\beta(4)+q_1+q_2,\\
 \widetilde a_5&=e_\beta(5)+q_2,\qquad
 \widetilde a_t=e_\beta(t)\quad(t\ge6),\\
 \widetilde\mu&=3-\ln6/\ln\beta+3q_0+4q_1+5q_2.
 \end{aligned}
\]
Set $\tau=345603/250000=1.382412$. Exact logarithm enclosures certify
that only indices $1,3$ have coefficient below $\tau$, while
$\widetilde a_2,\widetilde a_4,\widetilde a_5,e_\beta(6)>\tau$.
Monotonicity gives the bound for all later indices.
Put
\[
 \widetilde\chi=\tau-2\delta_\dagger
                [(\tau-1)+(\tau-\widetilde a_3)]
                -\delta_\dagger\widetilde\mu,\qquad
 \widetilde B=3+2\Phi_\beta(1)-\widetilde\chi.
\]
Lemma~\ref{lem:capped-rank-credit} now applies with exceptional
indices $E=\{1,3\}$, intercept $\widetilde\mu$, and
$\zeta=2/\ln\beta$. Its required condition
$\widetilde\chi+2/\ln\beta\ge\tau$ is certified below. The local
stage bound is therefore
$A\Delta+\widetilde B b+(I-J)/\ln2$.

\begin{theorem}[Mixed-prefix polynomial approximation]
\label{thm:poly-mixed-expected}
For general supplied binary unit-cost planar line queries and arbitrary
strictly positive rational priors, the oriented block algorithm with
threshold $\delta_\dagger$ has approximation factor
\[
 \boxed{\rho_{\mathrm{mix}}=
 A+\max\{\ln2/h(\delta_\dagger),
         \widetilde B\ln2/\sigma(\delta_\dagger)\}<5.283310.}
\]
It runs in polynomial time in the input and rational bit length.
\end{theorem}
\begin{proof}
Lemma~\ref{lem:stage-conversion} converts the local ledger using
\eqref{eq:poly-block-entropy}. Exact rational logarithm
bounds certify the coefficient comparisons, the zero-route condition, and
\[
 \begin{aligned}
 0.807181228455&<\widetilde\chi<0.807181228458,\\
 4.766748469916&<\widetilde B<4.766748469920,\\
 \max\{\ln2/h(\delta_\dagger),
         \widetilde B\ln2/\sigma(\delta_\dagger)\}
 &<1.691499244574.
 \end{aligned}
\]
Together with the previous bound on $A$, these give
$\rho_{\mathrm{mix}}<5.283309423415<5.283310$.
The convex-combination coefficients appear only in the proof. The
algorithm uses the same rational surrogate, supplied queries, and
polynomial shift enumeration as before, with its new rational balance
threshold. The complete certificate is included in
\texttt{poly-expected-next-check.py}.
\end{proof}

